%============================================================
% Chapter 9: Partitioning & In-Place Rearrangement
% STATUS: PLANNED  (generated by tools/outline.py; replace this file when the chapter is written)
%============================================================
\chapter{Partitioning \& In-Place Rearrangement}
\label{ch:partitioning}

\begin{tcolorbox}[enhanced, colback=white, colframe=secondaryblue, boxrule=1pt, arc=2mm]
\textbf{\large\textcolor{primaryblue}{Rearrange without a second array: partition schemes, the Dutch national flag, next permutation and interleavings.}}
\end{tcolorbox}

\chapterfacts{1100}{1900}{5}{10}{30}

\begin{plannedbox}[Planned Contents --- this chapter is not written yet]
\textbf{Part A --- Lessons}
\begin{enumerate}[leftmargin=*, label=\textbf{A\arabic*.}, itemsep=1pt]
    \item Partition invariants: Lomuto and Hoare, three-way partition
    \item Stable versus unstable rearrangement, and what extra memory buys
    \item Rearranging by sign, parity or index; interleaving two halves
    \item Next and previous permutation: the algorithm and its proof (ranking and the $k$-th permutation are the recursion book)
    \item Wiggle and zig-zag orders
\end{enumerate}
\textbf{Part B --- Patterns} (with candidate problems, checked when the chapter is written)
\begin{enumerate}[leftmargin=*, label=\textbf{P\arabic*.}, itemsep=1pt]
    \item \textbf{Three-way partition} --- Sort Colors {\small\color{darkgray}(LeetCode 75)}, Partition Array According to Given Pivot {\small\color{darkgray}(LeetCode 2161)}
    \item \textbf{Rearrange by a property} --- Rearrange Array Elements by Sign {\small\color{darkgray}(LeetCode 2149)}, Sort Array By Parity II {\small\color{darkgray}(LeetCode 922)}
    \item \textbf{Next permutation} --- Next Permutation {\small\color{darkgray}(LeetCode 31)}, Next Greater Element III {\small\color{darkgray}(LeetCode 556)}
    \item \textbf{Wiggle orders} --- Wiggle Sort II {\small\color{darkgray}(LeetCode 324)}
    \item \textbf{Rotations and block swaps}
\end{enumerate}
\end{plannedbox}
